\scanpage{39}
est déterminé par le complexe de
De Rham-Sullivan de $X_*$, et dans
le cas $k$ quelconque (et cas $2^\circ$ disons
-- $k$ principal est \unsure{le $1^\circ$|cas $1^\circ$|que $1^\circ$}) il est
déterminé par le complexe de
De Rham à p. d. \deleted{Enfin,}
\deleted{\unsure{il y aura|il y a}} (il est clair que $C^*$
avec les structures suppl. envisagées
\deleted{définit}\inserted{redétermine} le complexe de De Rham-Sullivan resp. le complexe de
De Rham à p. d.) Enfin, il
faut voir dans quelle mesure $C^*$
avec ses structures multiplicatives
à p. d. permet de reconstituer le
type d'homotopie du $X_*$ dont
il provient. (Enfin, bien sûr,
\unsure{faisceautiser|faisceauter} et relativiser -- !!!)
% archive sheet number: 38
